\section*{Chapter \ref{ch:iv1:order-statistics} --- Order Statistics}
\markboth{\protect\omTocMark{}{Answers: Order Statistics}}{\protect\omTocMark{}{Answers: Order Statistics}}

\begin{solution}{iq:iv1:order-statistics:1}
The largest averages $\frac{n}{n + 1}$, the smallest $\frac{1}{n + 1}$ and the range $\frac{n - 1}{n + 1}$. For $n = 10$, the largest is below $0.9$ with probability $0.9^{10} \approx 0.349$.

The $n$ points cut the interval into $n + 1$ pieces, and these pieces are exchangeable\index{uniform spacings}, so each has mean $\frac{1}{n + 1}$. The smallest is the first piece, the largest is everything except the last piece, and the range is everything except the two end pieces. The largest is below a value exactly when all $n$ are.

The case $n = 1$ checks all three: the largest and smallest both average $\tfrac12$, and the range is 0.

\iqwrong{Giving $1 - \tfrac1n$ for the largest. That is the point below which a single draw falls with chance $1 - 1/n$, not the mean of the largest. It is close for large $n$, and wrong for $n = 1$.}
\end{solution}

\begin{solution}{iq:iv1:order-statistics:2}
The density is $6x(1 - x)$, the mean is $\tfrac12$ and the variance is $1/20$. The average is better, with a variance of $1/36$.

The median is near $x$ when one point is near $x$, one below and one above. There are $3!$ ways of assigning those roles, which gives the density $6x(1 - x)$. In general, the $k$-th smallest of $n$ uniforms has a Beta$(k, n - k + 1)$ law\index{Beta law of order statistics}, here Beta$(2, 2)$. By symmetry the mean is $\tfrac12$, and $E[M^2] = \int_0^1 6x^3(1 - x)\,dx = 3/10$, so the variance is $3/10 - 1/4 = 1/20$.

The average of three has variance $\frac{1}{12 \cdot 3} = 1/36$, so the median is $1.8$ times as variable.

\iqwrong{Saying the median is always the more robust choice, so it must be the better one. Robustness protects against outliers, and a uniform has none. On a distribution with hard edges, the extremes are the most informative values.}
\end{solution}

\begin{solution}{iq:iv1:order-statistics:3}
The fastest belongs to desk B with probability $2/3$, and the three fastest all belong to it with probability $\binom{20}{3}/\binom{30}{3} = 57/203 \approx 0.281$. The distribution does not matter.

All 30 values come from one continuous distribution, so every ordering of them is equally likely\index{rank exchangeability}. The labels A and B are therefore a random assignment of 10 and 20 positions in the sorted list. The chance that the first three positions all go to B is the chain $\tfrac{20}{30} \cdot \tfrac{19}{29} \cdot \tfrac{18}{28}$.

This is what rank tests are based on. If the fastest three all came from one desk, you can compute that chance under ``same distribution'' with no model of latency at all.

\iqwrong{Cubing $2/3$, which gives about $0.296$. That treats the positions as independent, when they are drawn without replacement.}
\iqfollow{What is the test that this leads to?}{The Mann--Whitney rank test: it sums the ranks of one group and compares them with the distribution of all possible assignments.}
\end{solution}

\begin{solution}{iq:iv1:order-statistics:4}
The unbiased estimate is $60 \cdot (1 + \tfrac15) - 1 = 71$. Twice the mean minus one has a variance $7/3$ times as large, about $18^2$ against $12^2$ at $N = 71$.

The highest of $k$ serial numbers has $P(\max = x) = \binom{x - 1}{k - 1}/\binom{N}{k}$, and summing gives $E[\max] = k(N + 1)/(k + 1)$\index{German tank problem}. So $N = \frac{k + 1}{k} E[\max] - 1$, and replacing $E[\max]$ by what you saw gives an unbiased estimator.

There is a picture behind this. The $k$ numbers split $1, \dots, N$ into $k + 1$ gaps of the same average length. The numbers above the maximum are one gap, and the average of the $k$ gaps you can see is about $m/k$. So add that to $m$.

The sample mean has mean $(N + 1)/2$, so doubling it and taking off one is unbiased too, but the highest number carries all the information in the sample. Given the maximum, the other numbers are just uniform below it and say nothing more about $N$\index{sufficient statistic}. The mean rule ignores this, and with a sample $1, 2, 3, 4, 60$ it gives 27, which is below a serial number you have actually seen.

\iqwhy{Know which single number in a sample carries the information, and do not average it away. The same reasoning applies whenever a series has a hard edge, such as a position limit or a price floor.}
\end{solution}

\begin{solution}{iq:iv1:order-statistics:5}
The variances are $\theta^2/30$ and $\theta^2/120$ for $n = 10$, so the mean is 4 times as variable. The best multiple is $\frac{n + 2}{n + 1}\max X_i$, with mean squared error $\theta^2/(n + 1)^2$, here $\theta^2/121$ against $\theta^2/120$.

Take $\theta = 1$. The mean has variance $\frac{1}{12n}$, so $2\bar{X}$ has variance $\frac{1}{3n}$. The maximum has $E[M] = \frac{n}{n + 1}$ and $E[M^2] = \frac{n}{n + 2}$, so $\mathrm{Var}(M) = \frac{n}{(n + 1)^2(n + 2)}$, and the scaled maximum has variance $\frac{1}{n(n + 2)}$. The ratio is $(n + 2)/3$.

Note the rates: $1/n$ for the mean against $1/n^2$ for the maximum. An edge of a distribution is learned far faster than its centre\index{superefficiency at a boundary}.

For $cM$, the mean squared error is $c^2 \frac{n}{n + 2} - 2c \frac{n}{n + 1} + 1$, a parabola in $c$ with its minimum at $c = \frac{n + 2}{n + 1}$. This accepts a small bias for a lower variance, but the gain is tiny, $1/121$ against $1/120$\index{bias--variance trade-off}. The raw maximum, the maximum-likelihood estimate, has error $\frac{2}{(n + 1)(n + 2)} = 1/66$, nearly twice as much.

\iqwrong{Saying the maximum-likelihood estimate is the best because it usually is. Here it is biased low by construction, as it can never exceed the true $\theta$.}
\end{solution}

\begin{solution}{iq:iv1:order-statistics:6}
The second-highest of five bids averages $2/3$. Given that the highest is $0.8$, it averages $0.64$. The second-highest of four dice averages $5313/1296 \approx 4.10$.

From the gaps\index{uniform spacings}: the second-highest is everything except the top two of the six equal-mean pieces, so its mean is $4/6$.

Given the highest is $h$, the other four are independent uniforms on $[0, h]$, and the second-highest is the largest of them, with mean $\tfrac45 h$. At $h = 0.8$ that is $0.64$. As a check on the two answers, the highest has density $5h^4$, and averaging $\tfrac45 h$ over that gives $\tfrac45 \cdot \tfrac56 = \tfrac23$.

For the dice, use the tail sum\index{tail sum formula}. The second-highest is at least $k$ when at least two of the four dice are, so with $p = (7 - k)/6$,
\[
P(\text{second} \ge k) = 1 - (1 - p)^4 - 4p(1 - p)^3,
\]
and adding this for $k = 1, \dots, 6$ gives $5313/1296$.

\iqwrong{Using the uniform result for the dice, taking $1 + 5 \cdot \tfrac35 = 4$. A die is not a rescaled uniform, and ties make the answer slightly different.}
\iqwhy{In a second-price auction the winner pays the second-highest value, so this average is the seller's revenue. Knowing the top bid changes it, because the others are then bounded by it.}
\end{solution}

\begin{solution}{iq:iv1:order-statistics:7}
For a normal sample the variance of the median is $\pi/2 \approx 1.57$ times that of the mean. For the Laplace it is half that of the mean. So with fat enough tails the median wins.

The sample median has asymptotic variance $\frac{1}{4 n f(m)^2}$, with $f(m)$ the density at the true median\index{asymptotic variance of the median}. The reason is that around the median, the count of points below a level is binomial with variance $n/4$, and the slope of the expected count is $n f(m)$; dividing gives the error in position.

Standard normal: $f(0) = 1/\sqrt{2\pi}$, so the variance is $\frac{\pi}{2n}$, against $\frac1n$ for the mean.

Laplace: $f(0) = \tfrac12$, so the variance is $\frac1n$, while the mean has variance $\frac2n$.

\iqwrong{Saying the median is always less efficient because it ``throws data away''. It depends on the density at the centre compared with the tails. A peak at the centre and fat tails favour the median.}
\iqfollow{And at the extreme of fat tails?}{For a Cauchy sample the mean is no better than one draw, whatever $n$ is, while the median has variance $\frac{\pi^2}{4n}$. Real returns sit between the normal and these cases, which is why a trimmed mean or a median of a daily series is worth computing beside the mean.}
\end{solution}

\begin{solution}{iq:iv1:order-statistics:8}
The first fails after $1/5$ of a year on average, the third after $47/60$, the last after $137/60$. The time of the first failure and the time from it to the last are independent.

Five running servers fail at a total rate of 5, so the first failure takes $1/5$ on average. By lack of memory\index{exponential order statistics}, the remaining four then start afresh, with rate 4, and so on. So the sorted times are sums of independent gaps with means $\tfrac15, \tfrac14, \tfrac13, \tfrac12, 1$.

The third is $\tfrac15 + \tfrac14 + \tfrac13$, and the last is the sum of all five. The time from first to last is a sum of the other four gaps, so it is independent of the first, with mean $25/12$.

This independence is a special property of exponentials. In a uniform sample, a large minimum leaves little room, so the minimum and the range are strongly negatively correlated. A simulation finds about zero correlation for the exponentials, and below $-0.3$ for uniforms.

\iqwrong{Using one year for the last, the mean of a single lifetime. The last of five has to outlive the others, and that takes more than twice as long.}
\iqwhy{This is the standard picture for hardware failures with redundancy. With three replicas the time until the last is lost grows only slowly with redundancy, as a harmonic sum, unless you repair in between.}
\end{solution}

\begin{solution}{iq:iv1:order-statistics:9}
$P(N > n) = \frac{1}{n + 1}$. The smallest $n$ with $P(N \le n) \ge \frac12$ is 1, and the mean is infinite. The record events are independent.

$N > n$ means that the first draw is the largest of the first $n + 1$. All orderings are equally likely, so each of the $n + 1$ draws is the largest with the same chance\index{rank exchangeability}.

$P(N \le 1) = \tfrac12$ exactly, and $P(N \le 0) = 0$, so the smallest such $n$ is 1. The median is not unique, because $P(N \ge 2) = \tfrac12$ as well: every number from 1 to 2 is a median. The mean is the sum of the tail probabilities, $\sum_{n \ge 0} \frac{1}{n + 1}$, which diverges\index{tail sum formula}. If the first draw happens to be large, the wait is very long, and those rare cases dominate the average.

For the records: draw $k$ being a record is a statement about where the largest of the first $k$ lies. The relative order of the first 9 draws, in particular where the largest of the first 5 stands, is independent of whether draw 10 beats them all\index{independence of record events}. So the chance of both is $\tfrac15 \cdot \tfrac{1}{10} = \tfrac{1}{50}$, which a simulation confirms.

\iqwrong{Saying that the average wait must be finite because you nearly always finish after a few draws. The chance of a long wait falls too slowly, like $1/n$, for the sum to converge.}
\iqwhy{In a series of independent draws, a new all-time high gets rarer and rarer. A long run with no new high is expected under that model, and it says nothing on its own about a change.}
\end{solution}

\begin{solution}{iq:iv1:order-statistics:10}
Between the smallest and the largest: $1 - 2 \cdot 2^{-10} = 511/512 \approx 0.998$. Between the third smallest and the third largest: $912/1024 \approx 0.891$.

Each draw is below the true median with probability exactly $\tfrac12$, whatever the distribution\index{distribution-free interval}. So the number $B$ of draws below the median is binomial with $n = 10$ and $p = \tfrac12$.

The median exceeds the smallest draw unless $B = 0$, and it is below the largest unless $B = 10$, which gives the first answer. The median lies between the third smallest and third largest exactly when at least 3 draws are below it, and at least 3 above it, which means $3 \le B \le 7$:
\[
\frac{\binom{10}{3} + \binom{10}{4} + \binom{10}{5} + \binom{10}{6} + \binom{10}{7}}{2^{10}} = \frac{120 + 210 + 252 + 210 + 120}{1024} = \frac{912}{1024}.
\]
A simulation with exponential and normal draws gives the same two numbers.

\iqwrong{Using a normal interval on the sample, $\bar{X} \pm 2s/\sqrt{n}$. It needs assumptions you have not got. For a median, the binomial count gives an exact interval.}
\iqfollow{Does this work for the mean?}{No. A distribution can have the same median with any mean you like, because of a far-away small probability. No interval made from a finite sample can cover the mean under every distribution.}
\end{solution}

\begin{solution}{iq:iv1:order-statistics:11}
The best measured Sharpe ratio is about $1.08$, since the largest of 1000 standard normals averages about $3.24$. The expected best grows like $\sqrt{2 \ln n}$, very slowly.

Hand estimate: the largest of $n$ sits near the level exceeded with chance $1/n$\index{extreme value of normal samples}. Here that is $1/1000$, which is $3.09$ standard deviations, and using the median of the maximum, with tail $\ln 2/n$, it is $3.2$. Multiply by $\tfrac13$, and you get about $1.07$. The exact mean, from $\int x\, n\phi(x)\Phi(x)^{n - 1}dx$, is $3.24$, which a simulation matches.

The formula $\sqrt{2 \ln 1000} \approx 3.72$ has the right growth, but it is too high by half a standard deviation at this size. Going from 100 to 1000 strategies, or from 1000 to 10\,000, adds only about $0.6$ to $0.7$ standard deviation.

With one strategy there is no selection, and the expected best is 0.

\iqwrong{Saying that 1.08 over nine years is too high to be luck, because its $t$-statistic is above 3. That test is for one strategy chosen in advance. After 1000 tries you should expect that value when no strategy has any skill.}
\iqwhy{This is the first question to ask of a backtest: how many things were tried? A corrected hurdle needs the number of trials, and the true number includes the attempts nobody wrote down.}
\end{solution}

\begin{solution}{iq:iv1:order-statistics:12}
$P(R < r) = n r^{n - 1} - (n - 1) r^n$, and you need $n = 46$ draws, where the chance of falling short is $0.048$ (against $0.052$ for $n = 45$). For any continuous distribution, with 95\% confidence, at least 90\% of the distribution lies between the smallest and largest of 46 draws.

Take the range below $r$. Say the smallest point is the one at $x$, chosen in $n$ ways, and then the other $n - 1$ must lie in $[x, x + r]$. For $x \le 1 - r$ this has chance $r^{n - 1}$, and for $x$ above that, the chance is $(1 - x)^{n - 1}$. So
\[
P(R < r) = n\int_0^{1 - r} r^{n - 1}dx + n\int_{1 - r}^1 (1 - x)^{n - 1} dx = n r^{n - 1}(1 - r) + r^n,
\]
which is the formula. A check: for $n = 2$, it gives $2r - r^2$, which matches $|X - Y|$.

For a general distribution with cumulative function $F$, the values $F(X_i)$ are uniform\index{probability integral transform}. The mass between the smallest and largest sample is exactly the uniform range $F(\max) - F(\min)$, so the result holds with no assumption\index{tolerance interval}.

\iqwrong{Solving $0.9^n = 0.05$, which gives $n = 29$. That is right for one-sided cover, with only the maximum. Covering both ends takes more draws.}
\iqwhy{This gives a stress band with no model: 46 observed days bound 90\% of the distribution with 95\% confidence, whatever the returns look like, so long as the days are independent and from one distribution. They are often neither.}
\end{solution}

\begin{solution}{iq:iv1:order-statistics:13}
The variances are $1/132$ for the mean, $1/52$ for the median and $1/312$ for the midrange. The midrange wins, and its variance falls like $1/n^2$.

For the mean, the variance is $\frac{1}{12n}$. The median of $n = 2k + 1$ uniforms has a Beta$(k + 1, k + 1)$ law, with variance $\frac{1}{4(n + 2)}$, here $1/52$. For the midrange, the maximum and minimum have the same variance $\frac{n}{(n + 1)^2(n + 2)}$ and covariance $\frac{1}{(n + 1)^2(n + 2)}$, so
\[
\mathrm{Var}\Bigl(\frac{\max + \min}{2}\Bigr) = \frac{1}{2} \cdot \frac{n + 1}{(n + 1)^2(n + 2)} = \frac{1}{2(n + 1)(n + 2)}.
\]
A simulation matches all three within 1\%. For $n = 1$ all three coincide, at $1/12$, which is another check.

The midrange wins because the density has hard edges and is flat inside. The likelihood depends only on the two extremes, so given them, the rest of the points have no further information on $\theta$\index{sufficient statistic}. The median is worst, because it only uses local information, at a place with no special feature.

\iqwrong{Recommending the midrange in general. It is the worst choice with fat tails, and on a normal sample its variance is already above the mean's. A simulation at $n = 11$ shows it more than $1.5$ times higher. An estimator suits a shape, and there is no universally good one.}
\end{solution}

\begin{solution}{iq:iv1:order-statistics:14}
No two are within 10 ms with probability $0.91^{10} \approx 0.389$. The smallest gap averages $\frac{1}{n^2 - 1} = 1/99$ of a second, about $10.1$ ms.

The direct route is a shrinking argument\index{spacing argument}. Give the points in sorted order $x_1 < \dots < x_n$, all with gaps at least $d$. Take $y_i = x_i - (i - 1)d$. This is a one-to-one map from such sets to sorted sets of $n$ points in $[0, 1 - (n - 1)d]$, and it preserves volume. So the chance that all gaps are at least $d$ is
\[
\bigl(1 - (n - 1)d\bigr)^n,
\]
here $(1 - 9 \cdot 0.01)^{10}$. The formula gives 1 at $d = 0$ and 0 when $d = 1/(n - 1)$, when the points cannot fit any more.

The mean of the smallest gap is the integral of that tail, $\int_0^{1/(n - 1)} (1 - (n - 1)t)^n dt = \frac{1}{(n - 1)(n + 1)}$. So a chance of about $0.39$ at a distance of 10 ms sits beside a mean of about 10 ms, and this is a law with a long upper tail.

\iqwrong{Giving $(1 - 0.02)^{45} \approx 0.40$, by treating the 45 pairs as independent, each with chance $2d$ of being close. It is an approximation for rare events, and it comes out near the exact value, but here the pairs are not independent, and it would fail for a larger $d$.}
\iqwhy{Bursts occur at random. Ten orders in a second will come within 10 ms of each other more often than not, so a matching engine or a risk check has to be designed for that, even if the average gap is 100 ms.}
\end{solution}

\enlargethispage{-2\baselineskip}%
\begin{solution}{iq:iv1:order-statistics:15}
The largest of three averages $\frac{3}{2\sqrt{\pi}} \approx 0.846$, and the largest of four averages $\frac{3}{2\sqrt{\pi}}\bigl(1 + \frac{2}{\pi} \arcsin\frac13\bigr) \approx 1.029$. The method goes up to five (about $1.163$), and six would need the chance that four correlated normals are all positive, for which there is no simple formula.

Start with $E_n = n\int x\phi(x)\Phi(x)^{n - 1}dx$. Since $x\phi = -\phi'$, integrate by parts\index{Stein-type integration by parts}:
\[
E_n = n(n - 1) \int \phi(x)^2\, \Phi(x)^{n - 2} dx.
\]
Now, $\phi(x)^2 = \frac{1}{2\sqrt{\pi}}$ times the density of $W \sim N(0, \tfrac12)$. So
\[
E_n = \frac{n(n - 1)}{2\sqrt{\pi}}\, E\bigl[\Phi(W)^{n - 2}\bigr],
\]
and $\Phi(W)^{n - 2}$ is the chance that $n - 2$ new independent standard normals $A_i$ are all below $W$. The differences $A_i - W$ have variance $\tfrac32$, and any two of them have covariance $\tfrac12$, so correlation $\tfrac13$\index{orthant probability}. The chance that they are all negative is an orthant probability.

For $n = 3$: one variable, chance $\tfrac12$, so $E_3 = \frac{6}{2\sqrt{\pi}} \cdot \tfrac12$.

For $n = 4$: two variables, chance $\tfrac14 + \frac{1}{2\pi}\arcsin\frac13$, and multiplying by $\frac{12}{2\sqrt{\pi}}$ gives the stated formula.

For $n = 5$: three variables, chance $\tfrac18 + \frac{3}{4\pi}\arcsin\frac13$, so $E_5 = \frac{5}{4\sqrt{\pi}}\bigl(1 + \frac{6}{\pi}\arcsin\frac13\bigr)$. All three match a numerical integral to nine decimals.

A second route for $n = 3$ uses no integrals at all. The three pairwise distances add up to twice the range. Each has mean $2/\sqrt{\pi}$, so the range has mean $3/\sqrt{\pi}$, and the largest has half of that, because the smallest has the same size with the sign changed.

\iqwrong{Using $\sqrt{2 \ln n}$, which gives $1.48$ for $n = 3$. It is an asymptotic growth rate, and it is far from correct at small $n$. A quant should know the small values: $0.56$, $0.85$, $1.03$, $1.16$ for $n = 2$ to $5$.}
\end{solution}
